% Generated draft: requires mathematical and source review.
\documentclass[11pt]{article}
\usepackage[margin=25mm]{geometry}
\usepackage{fontspec}
\setmainfont{Latin Modern Roman}
\usepackage{amsmath,amssymb,mathtools,unicode-math}
\setmathfont{Latin Modern Math}
\usepackage{enumitem,needspace,xurl,hyperref}
\setlength{\parindent}{0pt}
\setlength{\parskip}{5pt}
\setlength{\emergencystretch}{3em}
\newcommand{\N}{\mathbb{N}}
\newcommand{\Z}{\mathbb{Z}}
\newcommand{\Q}{\mathbb{Q}}
\newcommand{\R}{\mathbb{R}}
\newcommand{\C}{\mathbb{C}}
\newcommand{\F}{\mathbb{F}}
\newcommand{\A}{\mathbb{A}}
\newcommand{\PP}{\mathbb{P}}
\newcommand{\E}{\mathbb{E}}
\newcommand{\eps}{\varepsilon}
\newcommand{\dd}{\,\mathrm{d}}
\begin{document}
\section*{UnsolvedMath ID: 30005160}
\textbf{Source problem code:} catalogue entry
\subsection{Definitions and mathematical statement}
Let \(d\geq 1\) be an integer and let \(\Gamma\subseteq\Z^d\). Write \(e(t)=\exp(2\pi i t)\) and \(n\cdot x=\sum_{j=1}^d n_jx_j\). For a finite positive real exponent \(p\), put
\[
\|f\|_{L^p([0,1]^d)}=\left(\int_{[0,1]^d}|f(x)|^p\,\dd x\right)^{1/p},
\]
where integration uses Lebesgue measure. For \(0<p<1\), this expression is a quasi-norm. The source does not explicitly specify whether its intended exponent range includes \(0<p<1\) or the endpoint \(p=\infty\); the formulas here describe finite positive exponents, subject to that unresolved scope qualification.

Using finitely supported coefficient families, say that \(\Gamma\) has the \emph{strict majorant property at exponent \(p\)} if, for every pair of real families \((a_n)_{n\in\Gamma}\) and \((A_n)_{n\in\Gamma}\) having finite support and satisfying \(|a_n|\leq A_n\) for every \(n\in\Gamma\), one has
\[
\left\|\sum_{n\in\Gamma}a_ne(n\cdot x)\right\|_{L^p([0,1]^d)}
\leq
\left\|\sum_{n\in\Gamma}A_ne(n\cdot x)\right\|_{L^p([0,1]^d)}.
\]
Thus \(A_n\geq0\). Despite the terminology, the defining inequality is non-strict. Finite support makes the sums unambiguous for arbitrary \(\Gamma\); it is an explicit formulation convention requiring confirmation, since the supplied excerpt does not state a support or convergence hypothesis for infinite \(\Gamma\).

\textbf{Classification question.} For each fixed admissible exponent \(p\notin\{2,4,6,\ldots\}\), characterize all sets \(\Gamma\subseteq\Z^d\) that satisfy the displayed majorant inequality. The target is a characterization at each fixed exponent, rather than solely a characterization of sets satisfying the property simultaneously at all exponents. The intended exponent range and the interpretation of infinite coefficient families remain unresolved as described above.

\textbf{Supplied contextual qualifications.} A set \(\Gamma\) is affinely independent if, whenever \(n_1,\ldots,n_k\) are distinct elements of \(\Gamma\) and real numbers \(c_1,\ldots,c_k\) satisfy \(\sum_jc_j=0\) and \(\sum_jc_jn_j=0\), all \(c_j\) vanish. The supplied dated assessment reports that the all-exponent classification is solved, with the strict majorant property at every exponent holding exactly for affinely independent frequency sets, and that broad counterexample families are known at non-even exponents. It expressly states that this theorem does not supply a complete characterization for each fixed non-even \(p\) and arbitrary \(\Gamma\). These are supplied literature assessments, not independently verified assertions here; the phrase ``every exponent'' inherits the source's unspecified exponent range.

The original excerpt also reports that, for every infinite \(\Gamma\subseteq\Z^d\), there are infinitely many positive integers \(m\) for which there exists a real coefficient family \((a_n)_{n\in\Gamma}\) such that, simultaneously for every \(p\in(2m,2m+2)\),
\[
\left\|\sum_{n\in\Gamma}|a_n|e(n\cdot x)\right\|_{L^p([0,1]^d)}
<
\left\|\sum_{n\in\Gamma}a_ne(n\cdot x)\right\|_{L^p([0,1]^d)}.
\]
The support or convergence conditions for these reported families are absent from the supplied excerpt. Finally, the excerpt contains the sentence ``In particular, this holds for every set \(\Gamma\subset\Z^d\) of cardinality at least \(d+2\).'' Its antecedent is missing, so this sentence cannot be assigned a precise mathematical conclusion from the supplied record.
\subsection{Short English statement}
For a fixed non-even exponent \(p\), determine which integer frequency sets \(\Gamma\subseteq\Z^d\) have the property that replacing real Fourier coefficients by nonnegative coefficientwise majorants never decreases the \(L^p\) size. The supplied assessment labels the problem partially solved: it distinguishes a reported all-exponent classification from the requested fixed-exponent characterization. The exponent range, support conditions for infinite sets, and one incomplete contextual assertion require review.
\subsection{Sources}
\begin{itemize}
\item Local UnsolvedMath registry: \texttt{data/unsolvedmath/problems.json}, permanent ID \(30005160\), exact source problem code \texttt{catalogue entry}, title \emph{Strict Majorant Properties at Fixed Noneven Exponents}; collection: Oberwolfach Reports Open Problems. The supplied proposed year is 2022; no proposer is identified.
\item Supplied report citation: \emph{Real Analysis, Harmonic Analysis and Applications} (2023), Oberwolfach Reports, DOI 10.4171/owr/2022/31, pp.\ 1743--1803. \url{https://doi.org/10.4171/owr/2022/31}. The background describes Oberwolfach Reports as documenting mathematical workshops and their open-problem discussions. The cited year 2023 is retained alongside the supplied proposed year 2022 and the DOI containing 2022; their relationship has not been independently checked.
\item Supplied literature reference, attributed to Gressman--Guo--Pierce--Roos--Yung, strict majorant property: \url{https://doi.org/10.1093/qmath/haac021}.
\item Supplied link identified as the same primary paper: \url{https://arxiv.org/abs/2106.12538}.
\item Supplied OWR report link: \url{https://ems.press/content/serial-article-files/46969}.
\item Literature assessment date: \texttt{2026-08-22}. The status ``partially solved,'' the reported all-exponent classification, and the counterexample assessment are supplied dated curation derived from the OWR record, not independently verified literature findings. The source expressly requests independent verification of time-sensitive status and cited claims.
\item The registry labels the statement \texttt{corrected\_verified} and says that only formula delimiters, notation, grammar, and typography were repaired, preserving mathematical content and scope. This is supplied source-verification metadata, not mathematical verification of this draft.
\end{itemize}
\end{document}
